% Standalone companion to the two-page AI4BCI submission.
% Original study, 13 October 2026; classical baseline added 14 October 2026.
% Compile twice with pdflatex.
\documentclass[11pt,a4paper]{article}
\usepackage[margin=0.75in]{geometry}
\usepackage{times,microtype,amsmath,amssymb,booktabs,enumitem,xcolor,hyperref}
\usepackage[numbers,sort&compress]{natbib}
\hypersetup{colorlinks=true,urlcolor=black,linkcolor=black,citecolor=black,
pdftitle={Companion Appendix: Recoverable Readout Mismatch in Smooth-to-Spiking EEG Decoding},
pdfauthor={Maximilian Nicholson}}
\setlength{\parindent}{0pt}
\setlength{\parskip}{5pt}
\setlist{nosep,leftmargin=*}
\renewcommand{\thetable}{A\arabic{table}}

\begin{document}
\begin{center}
{\Large\bfseries Companion Appendix:\\
Recoverable Readout Mismatch in\\Smooth-to-Spiking EEG Decoding\par}
\medskip
Maximilian Nicholson --- 14 October 2026
\end{center}

\textbf{Scope.} This document records the completed EEG study, supporting
analyses, and a subsequent implementation audit. The original primary
comparison remains calibrated SAST versus compute-matched ordinary training;
the readout analyses were planned secondary outcomes. The study was not
preregistered, and the narrative emphasis on recoverable readout mismatch
was chosen after inspecting the results. The source-moment diagnostic in
Appendix~\ref{sec:audit} was added after the main study.
A filter-bank common spatial pattern (FBCSP)-style reference was also added after the original run
and is reported separately in Appendix~\ref{sec:fbcsp}.

The \href{https://www.ai4bci.com/#cfp}{AI4BCI call} limits the submission to two
pages including references. This companion is kept separate and is not
assumed to be an accepted supplement. The main PDF is independently readable.

\appendix
\section{Data and neural dynamics}

\textbf{Dataset and partitions.} BCI Competition IV 2a, distributed as
BNCI2014\_001, contains nine participants and four motor-imagery classes
(left hand, right hand, feet, tongue) \citep{brunner}. Each of two sessions
has six runs of 48 trials, with 12 trials per class in each run. Only the
22 EEG channels enter the models; the three EOG channels are excluded.
All artifact-marked trials are retained. The artifact-free sensitivity
evaluates the same fitted models on the unmarked subset.

\begin{center}\small
\begin{tabular}{llll}
\toprule
Partition & Session / runs & Trials/person & Use \\
\midrule
Source training & 1 / 1--5 & 240 & Weights and preprocessing \\
Source validation & 1 / 6 & 48 & Settings and checkpoints \\
E1/E2 test & 2 / 1--6 & 288 & Conversion and probes \\
E3 calibration pool & 2 / 1 & 48 & Labeled session adaptation \\
E3 test & 2 / 2--6 & 240 & Fixed at every label budget \\
\bottomrule
\end{tabular}\end{center}

\textbf{Preprocessing.} MATLAB trial markers precede the cue by two seconds;
after converting to zero-based indexing, extraction adds 500 samples and
takes 1,000 samples at 250\,Hz. A fourth-order Butterworth 4--38\,Hz SOS filter
is applied forward/backward separately within each run, using odd padding
and SciPy's default pad length. Polyphase downsampling by two, with a Kaiser
window ($\beta=5$), produces $22\times500$ inputs covering 0--4\,s after cue
onset. Channel means and standard deviations are estimated from the 240
source-training epochs only. The forward/backward filter, padded temporal
convolutions, and across-time GroupNorm make this an offline evaluation,
distinct from the original competition's causal continuous-output scoring.

\textbf{Shared spiking-family architecture.} An eight-filter temporal
convolution (kernel $1\times31$) feeds spatial mixing across 22 electrodes
to 16 channels, average pooling by four, GroupNorm(1), and the first
leaky integrate-and-fire (LIF) block. A temporal convolution with kernel~5
maps 16 to 32 channels, followed by GroupNorm(1) and the second LIF block.
There are 125 time steps (32\,ms each); a linear 32-to-4 head reads their
mean activity. The model has 5,900 trainable parameters, no dropout, and no
batch running statistics. For each block,
\[
u_t=0.9u_{t-1}+I_t-\theta z_{t-1},\qquad
z_t=\begin{cases}
\tfrac12+\arctan(u_t-\theta)/\pi,&\text{smooth},\\
\mathbf{1}[u_t\geq\theta],&\text{hard}.
\end{cases}
\]
Initial voltage and activity are zero for each trial and each forward pass.
Raw conversion retains $\theta=1$ and all trained parameters. Changing
$z_t$ changes the subsequent reset feedback even though the recurrence
equation is unchanged. Direct-hard training uses hard forward activity
with the arctan derivative in the backward pass, including reset paths.

\clearpage
\section{Training, inference, and the original primary result}

Ordinary smooth training, same-batch SAM \citep{foret}, independent-batch SAST
\citep{nicholson}, and direct-hard training share architecture and initialization
for each participant/seed. SAM reuses one minibatch for ascent and descent;
SAST uses independently sampled minibatches. Original parameters are restored
before each base-optimizer update. The EEGNet-8,2 reference \citep{lawhern} is
an independent PyTorch implementation with a 32-sample temporal kernel,
depthwise spatial filters, separable temporal convolution, batch normalization,
ELU, pooling by 4/8, and dropout 0.5. Spatial and classifier max-norm limits
are 1 and 0.25. It has 2,164 parameters and uses PyTorch initialization and
the common optimizer, rather than reproducing the paper's full training
pipeline.

\textbf{Optimization and selection.} AdamW uses batches of 32 sampled with
replacement, 100 epochs of eight updates (200 epochs for matched ordinary),
and cosine decay to 1\% of the initial rate. Checkpoints are evaluated at 20
equally spaced positions in each method's budget, selected by native-forward
validation BA and then loss; all prescribed updates run. There is no gradient
clipping, mixed precision, or TF32. Each participant/method has four tuning
candidates on pilot seed 1000: learning rate $\{0.0003,0.001\}$ crossed with
radius $\{0.05,0.2\}$ at weight decay 0.0001 for SAM/SAST, or weight decay
$\{0.0001,0.001\}$ for the other methods. Ties between candidates use validation
loss and then candidate order. Main seeds are 0--4. All 216 pilot fits and
270 main fits completed before target-session evaluation.

\begin{table}[h]
\centering\small
\caption{Recorded main-fit compute on NVIDIA L4; means over 45 fits.
Time covers training and its validation, excluding data preparation, search,
and post-training analysis.}
\begin{tabular}{lrrr}
\toprule
Method & Optimizer updates & Gradient evaluations & Mean seconds \\
\midrule
% BEGIN GENERATED COST ROWS
Ordinary & 800 & 800 & 69.8 \\
Ordinary $2\times$ & 1600 & 1600 & 139.2 \\
SAM & 800 & 1600 & 139.4 \\
SAST & 800 & 1600 & 139.7 \\
Direct hard & 800 & 800 & 84.7 \\
EEGNet & 800 & 800 & 3.8 \\
% END GENERATED COST ROWS
\bottomrule
\end{tabular}
\end{table}

\textbf{Threshold calibration.} One threshold shared by both blocks is
selected from $\{0.25,0.5,0.75,1,1.25,1.5,2,3\}$ using source-validation BA.
Ties favor the closest value to 1, then the smaller value. Calibration
changes both firing threshold and reset magnitude. Direct-hard receives
the same calibration. The maximum grid value is selected in 17/45 ordinary,
18/45 matched, 12/45 SAM, and 14/45 SAST runs; conclusions concern this finite
grid, not optimal possible calibration.

\textbf{Inference.} BA is the unweighted mean of four class recalls.
Five seeds are averaged within each participant before computing means and
10,000 percentile bootstrap resamples (random seed 20261012). Contrasts are
paired within participant; the sampling unit is the participant ($n=9$).
Intervals are not corrected for the multiple secondary comparisons.
The primary effect, calibrated SAST minus calibrated matched ordinary,
is $-0.42$ pp, 95\% CI $[-2.88,2.67]$. The exact $2^9$ sign-flip sensitivity
test gives $p=0.8164$, conditional on sign exchangeability. Restricting
evaluation to artifact-free trials gives $-0.32$ pp $[-2.87,2.87]$.

\begin{center}\small
\begin{tabular}{lrrrrrrrrr}
\toprule
Participant & 1 & 2 & 3 & 4 & 5 & 6 & 7 & 8 & 9 \\
\midrule
% BEGIN GENERATED PRIMARY ROW
Difference (pp) & $+9.72$ & $+1.87$ & $-0.21$ & $+1.25$ & $-2.50$ & $-4.51$ & $-2.64$ & $-1.04$ & $-5.69$ \\
% END GENERATED PRIMARY ROW
\bottomrule
\end{tabular}\end{center}
The interval includes meaningful effects in either direction; it does not
establish equivalence. Neither this comparison nor the training algorithm
alone identifies a causal effect of parameter-space flatness.

\clearpage
\section{Conversion loss, readout recovery, and representation alignment}

\textbf{Matched probe construction.} Each mode and layer receives a separate
ridge classifier on time-averaged features. Source-training feature means and
standard deviations (floored at $10^{-6}$) define the scaler. A centered
least-squares fit to one-hot labels includes an unpenalized intercept, with
penalty $\alpha\in\{0.001,0.01,0.1,1,10,100,1000\}$ chosen by source-validation
BA; ties favor larger $\alpha$. All modes use the same 240 fitting and 48
validation trials. The 288 target trials are used only for evaluation.
The final-layer raw-versus-smooth probe contrast and its paired interval
were calculated during the final analysis from the planned probe outputs.

\begin{table}[h]
\centering\small\setlength{\tabcolsep}{4pt}
\caption{Within-checkpoint conversion loss and degenerate predictions.
Positive losses mean that smooth evaluation is more accurate. Collapse is
defined as predicting one class for all 288 target trials, out of 45 fits.}
\begin{tabular}{lcccc}
\toprule
& \multicolumn{2}{c}{Smooth minus spiking BA, pp [95\% CI]}
& \multicolumn{2}{c}{Single-class fits} \\
\cmidrule(lr){2-3}\cmidrule(lr){4-5}
Method & Raw & Calibrated & Raw & Calibrated \\
\midrule
% BEGIN GENERATED CONVERSION ROWS
Ordinary & 15.97 [8.86, 23.43] & 8.54 [4.96, 12.26] & 20 & 4 \\
Ordinary $2\times$ & 17.21 [10.22, 24.34] & 9.08 [4.77, 13.57] & 19 & 4 \\
SAM & 12.08 [5.25, 19.37] & 5.27 [1.99, 8.97] & 12 & 5 \\
SAST & 13.56 [7.13, 20.40] & 6.43 [3.29, 9.87] & 13 & 3 \\
% END GENERATED CONVERSION ROWS
\bottomrule
\end{tabular}
\end{table}

\begin{table}[h]
\centering\small\setlength{\tabcolsep}{4pt}
\caption{Final-layer linear probes: BA (\%) and paired effects (pp).
Recovery is raw-probe minus raw-original-head accuracy. All refitting uses
source-session labels only.}
\begin{tabular}{lrrrr}
\toprule
Method & Smooth probe & Raw probe & Recovery [95\% CI] & Raw $-$ smooth [95\% CI] \\
\midrule
% BEGIN GENERATED PROBE ROWS
Ordinary & 43.23 & 41.13 & 15.86 [9.23, 22.93] & -2.09 [-3.09, -1.08] \\
Ordinary $2\times$ & 43.67 & 42.37 & 17.04 [10.57, 23.46] & -1.30 [-2.35, -0.33] \\
SAM & 40.93 & 38.69 & 13.16 [6.28, 20.56] & -2.25 [-3.02, -1.55] \\
SAST & 41.91 & 40.53 & 14.62 [8.26, 21.37] & -1.37 [-2.03, -0.71] \\
% END GENERATED PROBE ROWS
\bottomrule
\end{tabular}
\end{table}

Refitting smooth features also improves accuracy: ordinary $+1.98$ pp,
matched ordinary $+1.13$ pp, SAM $+3.32$ pp, and SAST $+2.44$ pp.
This control limits how much of raw-spike recovery can be specifically
attributed to repairing conversion. The remaining negative raw-minus-smooth
probe contrasts support a residual loss of linear decodability under the
chosen readout family; they do not establish loss of all task information.
Calibrated-spike final-layer probe means are 40.27\%, 41.77\%, 38.94\%,
and 40.59\% in the same method order. Direct-hard raw features reach 44.30\%
with a ridge head; EEGNet features reach 54.88\%.

\textbf{Geometric similarity.} Linear centered kernel alignment (CKA)
\citep{kornblith} compares paired smooth and hard features on the target
trials, centered across trials. For centered feature matrices $X,Y$,
\[
\mathrm{CKA}(X,Y)=\frac{\|X^\top Y\|_F^2}
{\|X^\top X\|_F\,\|Y^\top Y\|_F}.
\]
A degenerate denominator is left undefined.
Final-layer CKA can remain high despite collapsed original-head predictions.

\begin{center}\small\setlength{\tabcolsep}{5pt}
\begin{tabular}{lrrrr}
\toprule
& \multicolumn{2}{c}{Layer 1 CKA} & \multicolumn{2}{c}{Layer 2 CKA} \\
\cmidrule(lr){2-3}\cmidrule(lr){4-5}
Method & Raw & Calibrated & Raw & Calibrated \\
\midrule
% BEGIN GENERATED CKA ROWS
Ordinary & 0.992 & 0.981 & 0.903 & 0.890 \\
Ordinary $2\times$ & 0.994 & 0.985 & 0.876 & 0.879 \\
SAM & 0.947 & 0.935 & 0.801 & 0.823 \\
SAST & 0.986 & 0.983 & 0.900 & 0.914 \\
% END GENERATED CKA ROWS
\bottomrule
\end{tabular}\end{center}
SAST final-layer raw CKA is 0.900 [0.865, 0.928], and calibrated CKA is
0.914 [0.883, 0.935]. These are descriptive geometric measures; high CKA
does not establish preserved decisions, identical dynamics, or equivalence.

\clearpage
\section{Limited-label target-session adaptation}

E3 is a planned secondary experiment with a different test partition from
E1/E2. All of target run~1 is reserved for calibration at every budget,
including zero. Nested calibration sets take the first $k=0,4,8,12$ trials
per class in recording order, giving 0, 16, 32, or 48 labels. Target runs~2--6
(240 trials) form the fixed test set. No unused trial from the calibration
run enters the test set.

\textbf{Readout only.} The encoder and source-selected threshold are frozen.
The ridge head is refitted on source-training plus calibration trials with
equal trial weights; feature scaling and $\alpha$ remain fixed from the
source-only fit. The zero-label baseline is already a source-fitted ridge
head, rather than the original trained classifier. Accordingly, gains below
measure the extra benefit of target labels after source-only refitting.

\begin{table}[h]
\centering\small\setlength{\tabcolsep}{5pt}
\caption{Readout adaptation on a fixed target test set. Entries are mean BA
(\%). The final column is the paired 48-minus-0-label gain, pp [95\% CI].
The spiking modes use source-selected thresholds.}
\begin{tabular}{llrrrrr}
\toprule
Method & Mode & 0 & 16 & 32 & 48 & Gain [95\% CI] \\
\midrule
% BEGIN GENERATED ADAPTATION ROWS
Ordinary & Smooth & 43.08 & 43.51 & 43.72 & 44.33 & 1.25 [0.45, 2.12] \\
Ordinary & Spikes & 40.19 & 40.49 & 41.09 & 41.49 & 1.30 [0.42, 2.37] \\
Ordinary $2\times$ & Smooth & 43.38 & 43.75 & 44.19 & 44.19 & 0.81 [-0.09, 1.94] \\
Ordinary $2\times$ & Spikes & 41.54 & 41.76 & 42.30 & 42.43 & 0.89 [0.17, 1.62] \\
SAM & Smooth & 40.88 & 41.23 & 41.45 & 41.68 & 0.80 [-0.07, 1.60] \\
SAM & Spikes & 38.90 & 39.11 & 39.38 & 39.35 & 0.45 [-0.20, 1.11] \\
SAST & Smooth & 41.96 & 42.45 & 42.75 & 43.11 & 1.15 [0.48, 1.81] \\
SAST & Spikes & 40.47 & 40.97 & 41.06 & 41.46 & 0.99 [0.43, 1.54] \\
Direct hard & Spikes & 44.44 & 44.90 & 45.11 & 45.24 & 0.81 [0.32, 1.28] \\
EEGNet & Native & 53.88 & 54.15 & 54.63 & 55.50 & 1.62 [0.94, 2.42] \\
% END GENERATED ADAPTATION ROWS
\bottomrule
\end{tabular}
\end{table}

\textbf{Threshold only.} All parameters, including the original head, are
fixed. At zero labels the source-selected threshold is used; at positive
budgets the same eight-candidate grid is selected by calibration-trial BA,
with ties closest to the source threshold and then toward the smaller value.
These zero-label baselines differ from the readout branch.

\begin{table}[h]
\centering\small
\caption{Threshold-only adaptation, mean BA (\%), with the same fixed test
set and label budgets. The final column is the paired 48-minus-0-label gain.}
\begin{tabular}{lrrrrr}
\toprule
Method & 0 & 16 & 32 & 48 & Gain, pp [95\% CI] \\
\midrule
% BEGIN GENERATED THRESHOLD ROWS
Ordinary & 32.72 & 32.30 & 32.06 & 32.31 & -0.42 [-1.94, 0.60] \\
Ordinary $2\times$ & 33.44 & 33.31 & 32.67 & 33.45 & 0.01 [-0.42, 0.46] \\
SAM & 32.18 & 32.34 & 32.48 & 33.00 & 0.82 [-0.14, 2.04] \\
SAST & 33.02 & 32.46 & 32.78 & 33.26 & 0.24 [-0.74, 1.14] \\
Direct hard & 42.97 & 40.13 & 41.22 & 42.44 & -0.54 [-1.30, 0.46] \\
% END GENERATED THRESHOLD ROWS
\bottomrule
\end{tabular}
\end{table}

For calibrated SAST, 48 target labels add $0.99$ pp $[0.43,1.54]$ after
source-only ridge refitting, compared with $0.24$ pp $[-0.74,1.14]$ for
threshold-only adaptation. These increments are smaller than the source-only
readout recovery in E2, but cannot be directly subtracted from E2 because
the target test partitions differ. Calibration budgets, interventions,
and secondary contrasts were not optimized on the held-out test trials.
The approach assumes retained access to the source training data.

\clearpage
\section{Implementation audit and exploratory source alignment}
\label{sec:audit}

\textbf{Audit coverage.} The archived engine, local source, and engine embedded
in the notebook have identical SHA-256 hashes. All 675 E1 participant/seed/mode
BA values and predicted-class counts were recomputed from archived confusion
matrices and agreed to numerical precision. For participant~1, seed~0,
three checkpoints (SAST, matched ordinary, direct-hard) were replayed on CPU.
All predictions matched across eight applicable mode evaluations; maximum
absolute feature differences were below $6.6\times10^{-7}$.

The sampled audit also reproduced training-only input normalization and
all 24 source-validation threshold-grid scores, checked trial-state resets
and inference-batch invariance, and confirmed that conversion left weights
unchanged. An independent scalar recurrence agreed with the implemented
smooth and hard dynamics. This audit found no discrepancy explaining the
observed collapse within its tested scope. It was not an exhaustive replay
of all 270 main checkpoints and did not retrain models.

\textbf{Exploratory alignment.} To test whether simple feature offsets and
scales could explain part of the mismatch, the original classifier was held
fixed for the SAST and matched-ordinary checkpoints above. Let $\mu_r,\sigma_r$
and $\mu_s,\sigma_s$ be per-feature source-training moments under raw and
smooth dynamics. Mean-only correction uses
$z'_r=z_r-\mu_r+\mu_s$. Mean-and-scale correction uses
\[
z'_r=(z_r-\mu_r)\odot
\frac{\sigma_s}{\max(\sigma_r,10^{-6})}+\mu_s .
\]
Moments use the 240 source-training trials without their labels; no target
labels or target moments enter the transform. These diagnostics were added
after inspecting study outcomes and evaluated on a convenience sample of
one participant and one seed, rather than across the full benchmark.

\begin{table}[h]
\centering\small
\caption{Exploratory participant-1/seed-0 BA (\%) with the original head.
These two checkpoints are illustrative and are not an aggregate result.}
\begin{tabular}{lrrrr}
\toprule
Method & Raw & Mean only & Mean and scale & Smooth reference \\
\midrule
SAST & 25.35 & 38.89 & 47.57 & 53.47 \\
Matched ordinary & 25.00 & 38.54 & 45.49 & 55.56 \\
\bottomrule
\end{tabular}
\end{table}

For this SAST checkpoint, the mean across-feature trial standard deviation
falls from 0.257 in smooth mode to 0.088 in raw mode; 41.2\% of raw
time-averaged feature entries are zero. Spiking activity is therefore not
universally absent. The alignment response is consistent with a readout
mismatch involving feature moments, but does not establish a general repair
method or isolate a causal mechanism. A full participant/seed evaluation
of fixed source-moment correction remains unrun.

\textbf{Interpretive limits.} The evidence concerns one dataset, one
convolutional/LIF architecture, source-selected settings, and an offline
participant-specific session shift. Probe success does not prove identical
representations; residual probe differences do not identify all information
loss. Artifact-free evaluation cannot rule out every ocular or muscle
contribution. EEGNet and the FBCSP-style reference remain more accurate
than the converted models under these implementations.
No hardware energy, real-time latency, cross-participant transfer, or
clinical efficacy was measured.

\clearpage
\section{Post-study classical decoding reference}
\label{sec:fbcsp}

\textbf{Scope and method.} An additional FBCSP-style reference follows the
filter-bank/common spatial pattern (CSP) structure of Ang et al.\ \citep{ang}, adapted to the existing
offline front end. The source-training runs~1--5, source-validation run~6,
all-session-2 test set, cue epochs, channel scaling, and trial retention
match E1/E2. Nine fourth-order Butterworth bands are applied forward/backward
within each 125\,Hz epoch: 4--8, 8--12, 12--16, 16--20, 20--24, 24--28,
28--32, 32--36, and 34--38\,Hz. The last band overlaps the preceding band
to stay within 4--38\,Hz.

Each band has four one-vs-rest CSP models, with four extreme components
per model, covariance shrinkage 0.1, full rank, and no trace normalization.
The 144 log-mean-square-power features are standardized with source-training
statistics. Source-training mutual information (three neighbors; random
seed 20261012) selects 8, 16, 32, or 64 features. Multiclass linear discriminant
analysis (LDA) uses the
least-squares solver and automatic covariance shrinkage. Validation BA
selects a candidate, with ties resolved by validation log loss and then
candidate order. The selected model is not refitted on validation trials.
All nine source choices are locked before session~2 evaluation. This is one
deterministic baseline per participant; EEGNet retains its original five-seed
participant mean. The original FBCSP competition method used causal
Chebyshev filtering and a Naive Bayesian Parzen Window classifier, so this
reference is an adaptation rather than an exact reproduction.

\textbf{Results.} FBCSP-style BA is 61.81\% [52.93, 70.06], compared with
54.14\% [43.83, 64.21] for EEGNet. The paired participant-level difference is
$+7.66$ pp [3.67, 11.91]. Intervals use the same 10,000 participant bootstrap
resamples and seed as the original analysis. This comparison was added
after the original study; its intervals are descriptive, not corrected for
multiple comparisons, and do not replace the primary SAST contrast.

\begin{center}\small
\begin{tabular}{rrrrr}
\toprule
Participant & Features & FBCSP BA (\%) & EEGNet BA (\%) & Difference (pp) \\
\midrule
% BEGIN GENERATED FBCSP PARTICIPANT ROWS
1 & 32 & 75.00 & 67.78 & +7.22 \\
2 & 32 & 44.10 & 42.57 & +1.53 \\
3 & 32 & 72.57 & 70.49 & +2.08 \\
4 & 64 & 59.72 & 41.25 & +18.47 \\
5 & 8 & 45.14 & 31.11 & +14.03 \\
6 & 64 & 42.71 & 34.72 & +7.99 \\
7 & 64 & 73.26 & 58.06 & +15.21 \\
8 & 64 & 68.75 & 66.94 & +1.81 \\
9 & 32 & 75.00 & 74.38 & +0.63 \\
% END GENERATED FBCSP PARTICIPANT ROWS
\bottomrule
\end{tabular}
\end{center}

\textbf{Provenance.} The completed baseline run is
\texttt{FBCSP\_LDA\_fb436529d768932f}. It used Python 3.13.16,
NumPy 2.1.3, SciPy 1.16.3, pandas 2.2.3, MNE 1.10.2, and scikit-learn 1.7.2.
The engine SHA-256 is:
\begin{center}\small\ttfamily
4c694da2ab6d74d1dcb8a7428b5d240b\\
f6572c2a71aea9d91f45403645722cd3
\end{center}
The archived engine matches the locally tested implementation. An independent
export audit recomputed all nine BA scores and paired intervals from 2,592
saved test predictions and checked model, selection, reference, and split
records. Model files were hash-checked; the nine fits were not independently
retrained for this audit.

\subsection{Post-study EEGNet source-validation audit}

The original EEGNet reference was audited without using target labels or changing
the training procedure. All 45 main fits and 36 tuning fits completed 800 updates
and 20 source-validation evaluations. Recomputed checkpoint choices and all nine
optimizer-grid winners agree with the archive. Finite weights, spatial/head
max-norm constraints, and BatchNorm update counts also pass. Across main fits,
the median selected epoch is 80; 20/45 select epoch 90 or later, including 16/45
at epoch 100. Minimum validation loss occurs at epoch 100 in 20/45 fits, and
13/45 improve their best validation accuracy in the final 25 epochs.
These histories do not establish convergence or that longer training would
improve target-session performance. EEGNet is a reference under the common
study budget and four-candidate source search, not a tuned benchmark ceiling.
The audit script, all source learning curves, and input hashes accompany the release.

\clearpage
\section{Post-study second-architecture readout check}

This exploratory check asks whether readout recovery also occurs in a structurally
different spiking encoder. It uses the same nine participants, source/target session
split, artifact retention, and offline preprocessing. It is not an independent
dataset replication and does not repeat the SAST comparison.

\textbf{Architecture and fixed protocol.} A 63-sample temporal convolution has
eight filters shared across electrodes. Depthwise spatial filtering learns two
electrode combinations per temporal filter, producing 16 channels. Average pooling
by eight and GroupNorm(1) precede one LIF block (leak 0.9, threshold 1, zero initial
state, subtractive previous-activity reset). Four temporal-bin averages form a
64-dimensional feature vector and a linear four-class head. The encoder has
1,164 trainable parameters. Linear spatial filtering is evaluated before temporal
filtering to reduce intermediates; a reference-value and gradient check verifies
the equivalent filter ordering. This differs from the original encoder's two LIF
blocks, full spatial mixing, and one temporal average per final channel.

Only ordinary smooth training is used. The original 800-update AdamW budget,
batch sampling, cosine schedule, four learning-rate/weight-decay candidates,
pilot seed 1000, and best-checkpoint rule are retained. Five main seeds per
participant use the source-selected setting: 36 tuning fits and 45 main fits.
Matched ridge heads are fitted separately to smooth and raw-spike features using
the original source-training data and seven-value validation penalty grid.
All participants' source fits and readout selections are locked before loading
target sessions; file hashes are checked again after evaluation. Encoders remain
frozen. No target labels tune the architecture, optimizer, checkpoint, or readout.

The design was fixed before this check's target evaluation, after the original
study and its classical reference had been inspected. It was not preregistered.
Seeds are averaged within participant before the same 10,000-resample participant
bootstrap. Smooth-refit recovery is retained as the control for generic refitting
benefits; the excess raw recovery is a paired difference of recovery gains.

% BEGIN GENERATED ARCHITECTURE SUMMARY
Raw original-head BA is 27.44 [26.25, 28.78]\%; source-only raw-spike ridge BA is 37.55 [32.59, 42.61]\%, giving a paired gain of 10.11 [5.86, 14.29] pp. Smooth original-head BA is 36.44 [31.60, 41.57]\% and smooth-ridge BA is 39.46 [34.98, 43.85]\%, with gain 3.02 [0.70, 5.42] pp. The raw-minus-smooth matched-ridge contrast is -1.91 [-3.14, -0.61] pp; excess raw recovery over smooth recovery is 7.09 [2.01, 12.21] pp. Intervals are paired participant-bootstrap intervals. These estimates describe this architecture check; they do not add independent participants or establish generality across datasets.
% END GENERATED ARCHITECTURE SUMMARY

\begin{table}[h]
\centering\small
\caption{Second architecture: session~2 balanced accuracy (\%), averaged over five
seeds within each participant. Original and ridge heads share the same frozen encoder.}
\begin{tabular}{r r r r r}
\toprule
Participant & Smooth original & Smooth ridge & Raw original & Raw ridge \\
\midrule
% BEGIN GENERATED ARCHITECTURE ROWS
1 & 40.56 & 43.19 & 29.38 & 43.40 \\
2 & 26.11 & 28.89 & 26.11 & 25.07 \\
3 & 51.67 & 51.25 & 31.67 & 51.88 \\
4 & 30.49 & 40.83 & 27.85 & 36.25 \\
5 & 26.11 & 32.78 & 24.79 & 28.96 \\
6 & 34.86 & 31.46 & 27.08 & 32.29 \\
7 & 35.90 & 38.89 & 27.99 & 35.90 \\
8 & 41.18 & 43.82 & 25.62 & 41.74 \\
9 & 41.11 & 44.03 & 26.46 & 42.43 \\
% END GENERATED ARCHITECTURE ROWS
\bottomrule
\end{tabular}
\end{table}

\clearpage
\section{Reproducibility, provenance, and declarations}

The original run, \texttt{BNCI2014\_001\_ccebe0c6749374ab}, completed on
13 October 2026 using NVIDIA L4, Python 3.13.15, and PyTorch 2.11.0+cu130.
The CPU audit and second architecture use Python 3.12.12 and PyTorch 2.14.0;
full dependency versions are archived with each run.
The experiment-engine SHA-256 is:
\begin{center}\small\ttfamily
42ba021e26724e3e0d18ec77d1df765cc4380\\
92cf487c2eb1b94b8d5a56aacde
\end{center}
% Concatenate the two hash lines without whitespace.

Public reproduction package:\\
\url{https://github.com/PluvioXO/ai4bci2026-readout}.\\
Full run archives are versioned release assets. The package contains:
\begin{itemize}
\item \texttt{ai4bci\_all\_experiments.ipynb}: the standalone original experiment
runner with its exact archived engine.
\item \texttt{results/} and \texttt{build\_paper\_results.py}: per-seed outputs,
training costs, and regeneration of paper figures, tables, and paired intervals.
\item \texttt{conversion\_audit/}: numerical checks, sampled forward replay,
and explicitly exploratory alignment results.
\item \texttt{ai4bci\_fbcsp\_baseline.ipynb} and \texttt{results/fbcsp\_baseline/}:
the classical runner, source selections, models, predictions, and export audit.
\item \texttt{ai4bci\_readout\_replication.ipynb} and
\texttt{results/architecture\_replication/}: the standalone second-architecture
runner and its participant/seed results, intervals, fixed protocol, and export audit.
\item \texttt{code/}: both neural engines, source-history and prediction audits,
and preparation of the three-checkpoint conversion replay.
\end{itemize}
Full archives retain source configuration, input checksums, splits, histories,
selected weights, features, predictions, and calibration identifiers. The main
paper compiles independently with standard LaTeX and PGFPlots. Dataset recordings
and prepared EEG epochs are not redistributed.

\textbf{Relationship to other work.} SAST was introduced in the author's
event-based and on-sensor spiking-network preprints \citep{nicholson,nicholsonOnSensor}. The present paper applies SAST to EEG motor imagery and
adds a planned source-only readout analysis on frozen spiking features. The
related UniReps manuscript, \emph{On Representational Fidelity Across Neural
Dynamics: Does Flatness Preserve Continuous-to-Spiking Computation?}, is
submitted but unpublished and addresses a separate question.

\textbf{Data and ethics.} This is secondary analysis of publicly distributed
EEG recordings; no new participants were recruited and no additional
recordings were collected. The dataset is obtained from its original
distribution and is cited in the main paper. No institutional approval or
exemption number is asserted here. The applicable institutional secondary-analysis
determination remains unconfirmed.

\begin{thebibliography}{7}
\setlength{\itemsep}{0pt}
\setlength{\parskip}{0pt}
\bibitem{brunner}
C.~Brunner, R.~Leeb, G.~R.~M\"uller-Putz, A.~Schl\"ogl, and G.~Pfurtscheller.
\href{https://www.bbci.de/competition/iv/desc_2a.pdf}{BCI Competition 2008---Graz
data set A}. Dataset description, 2008.

\bibitem{foret}
P.~Foret, A.~Kleiner, H.~Mobahi, and B.~Neyshabur.
\href{https://arxiv.org/abs/2010.01412}{Sharpness-Aware Minimization for
Efficiently Improving Generalization}. ICLR, 2021.

\bibitem{nicholson}
M.~Nicholson. \href{https://arxiv.org/abs/2603.18039}{Sharpness Aware Surrogate
Training for Spiking Neural Networks}. arXiv:2603.18039, 2026.

\bibitem{nicholsonOnSensor}
M.~Nicholson. \href{https://arxiv.org/abs/2604.09696}{Sharpness-Aware
Surrogate Training for On-Sensor Spiking Neural Networks}.
arXiv:2604.09696, 2026.

\bibitem{lawhern}
V.~J.~Lawhern et~al.
\href{https://doi.org/10.1088/1741-2552/aace8c}{EEGNet: A compact convolutional
neural network for EEG-based brain--computer interfaces}.
J. Neural Eng., 15:056013, 2018.

\bibitem{kornblith}
S.~Kornblith, M.~Norouzi, H.~Lee, and G.~Hinton.
\href{https://proceedings.mlr.press/v97/kornblith19a.html}{Similarity of Neural
Network Representations Revisited}. ICML, pp.~3519--3529, 2019.

\bibitem{ang}
K.~K.~Ang, Z.~Y.~Chin, C.~Wang, C.~Guan, and H.~Zhang.
\href{https://doi.org/10.3389/fnins.2012.00039}{Filter Bank Common Spatial Pattern
Algorithm on BCI Competition IV Datasets 2a and 2b}. Front. Neurosci., 6:39, 2012.
\end{thebibliography}
\end{document}
